\section{Introducción a la estadística descriptiva}

La \emph{estadística} es la ciencia que se encarga de recopilar, organizar, analizar e interpretar datos para tomar decisiones informadas. Se divide en dos grandes ramas:

\begin{itemize}
	\item \textbf{Estadística descriptiva:} Se enfoca en resumir y describir las características de un conjunto de datos mediante medidas numéricas y representaciones gráficas.
	\item \textbf{Estadística inferencial:} Utiliza una muestra de datos para hacer generalizaciones y predicciones sobre una población más grande.
\end{itemize}

En este capítulo nos enfocaremos en la estadística descriptiva, que es la base para comprender conceptos más avanzados de probabilidad e inferencia estadística.

\subsection{¿Por qué es importante la estadística descriptiva?}

En la era de los datos, la capacidad de resumir y visualizar información es fundamental. La estadística descriptiva nos permite:

\begin{itemize}
	\item \textbf{Resumir grandes volúmenes de datos} en unas pocas medidas representativas.
	\item \textbf{Identificar patrones y tendencias} en los datos.
	\item \textbf{Comparar diferentes conjuntos de datos} de manera objetiva.
	\item \textbf{Detectar valores atípicos} (outliers) que pueden indicar errores o fenómenos interesantes.
	\item \textbf{Comunicar resultados} de forma clara y concisa.
\end{itemize}

\subsection{Conceptos fundamentales}

Antes de estudiar las medidas descriptivas, es importante distinguir entre:

\begin{definicion}[Población y muestra]
	\begin{itemize}
		\item \textbf{Población:} Es el conjunto completo de elementos o individuos que comparten una característica común y sobre los cuales queremos obtener información.
		\item \textbf{Muestra:} Es un subconjunto representativo de la población que se selecciona para el estudio.
	\end{itemize}
\end{definicion}

\begin{definicion}[Parámetro y estadístico]
	\begin{itemize}
		\item \textbf{Parámetro:} Es una medida numérica que describe una característica de la \emph{población} (por ejemplo, la media poblacional $\mu$).
		\item \textbf{Estadístico:} Es una medida numérica calculada a partir de una \emph{muestra} (por ejemplo, la media muestral $\bar{X}$).
	\end{itemize}
\end{definicion}

\begin{observacion}
	En la práctica, rara vez podemos estudiar toda la población, por lo que utilizamos estadísticos calculados de muestras para estimar los parámetros poblacionales.
\end{observacion}

\subsection{Tipos de medidas descriptivas}

Las medidas descriptivas se clasifican en:

\begin{enumerate}
	\item \textbf{Medidas de tendencia central:} Indican el valor alrededor del cual se concentran los datos (media, mediana, moda).
	\item \textbf{Medidas de dispersión:} Indican qué tan dispersos están los datos (rango, desviación estándar, varianza).
	\item \textbf{Medidas de posición:} Indican la posición relativa de un valor dentro del conjunto de datos (cuartiles, percentiles).
	\item \textbf{Medidas de forma:} Describen la forma de la distribución de los datos (asimetría, curtosis).
\end{enumerate}

En las siguientes secciones estudiaremos las medidas de tendencia central y dispersión, que son las más utilizadas en la práctica.
